\IfFileExists{stacks-project.cls}{%
\documentclass{stacks-project}
}{%
\documentclass{amsart}
}
\usepackage{amssymb}

% For dealing with references we use the comment environment
\usepackage{verbatim}
\newenvironment{reference}{\comment}{\endcomment}
%\newenvironment{reference}{}{}
\newenvironment{slogan}{\comment}{\endcomment}
\newenvironment{history}{\comment}{\endcomment}

% For commutative diagrams we use Xy-pic
\usepackage[all]{xy}

% We use 2cell for 2-commutative diagrams.
\xyoption{2cell}
\UseAllTwocells

% We use multicol for the list of chapters between chapters
\usepackage{multicol}

% This is generally recommended for better output
\usepackage{lmodern}
\usepackage[T1]{fontenc}
\usepackage{textalpha}

% For cross-file-references
\usepackage{xr-hyper}
\newcommand{\maybeexternaldocument}[2]{%
  \IfFileExists{#2.aux}{%
    \externaldocument[#1]{#2}%
  }{%
    \PackageWarning{stack}{Missing #2.aux; external references with prefix #1 unavailable}%
  }%
}
\let\externaldocumentorig\externaldocument
\renewcommand{\externaldocument}[2][]{%
  \IfFileExists{#2.aux}{%
    \externaldocumentorig[#1]{#2}%
  }{%
    \PackageWarning{stack}{Missing #2.aux; external references with prefix #1 unavailable}%
  }%
}

% Package for hypertext links:
\usepackage{hyperref}

% For any local file, say "hello.tex" you want to link to please
% use \externaldocument[hello-]{hello}
\externaldocument[introduction-]{introduction}
\externaldocument[conventions-]{conventions}
\externaldocument[sets-]{sets}
\externaldocument[categories-]{categories}
\externaldocument[topology-]{topology}
\externaldocument[sheaves-]{sheaves}
\externaldocument[sites-]{sites}
\externaldocument[stacks-]{stacks}
\externaldocument[fields-]{fields}
\externaldocument[algebra-]{algebra}
\externaldocument[brauer-]{brauer}
\externaldocument[homology-]{homology}
\externaldocument[derived-]{derived}
\externaldocument[simplicial-]{simplicial}
\externaldocument[more-algebra-]{more-algebra}
\externaldocument[smoothing-]{smoothing}
\externaldocument[modules-]{modules}
\externaldocument[sites-modules-]{sites-modules}
\externaldocument[injectives-]{injectives}
\externaldocument[cohomology-]{cohomology}
\externaldocument[sites-cohomology-]{sites-cohomology}
\externaldocument[dga-]{dga}
\externaldocument[dpa-]{dpa}
\externaldocument[sdga-]{sdga}
\externaldocument[hypercovering-]{hypercovering}
\externaldocument[schemes-]{schemes}
\externaldocument[constructions-]{constructions}
\externaldocument[properties-]{properties}
\externaldocument[morphisms-]{morphisms}
\externaldocument[coherent-]{coherent}
\externaldocument[divisors-]{divisors}
\externaldocument[limits-]{limits}
\externaldocument[varieties-]{varieties}
\externaldocument[topologies-]{topologies}
\externaldocument[descent-]{descent}
\externaldocument[perfect-]{perfect}
\externaldocument[more-morphisms-]{more-morphisms}
\externaldocument[flat-]{flat}
\externaldocument[groupoids-]{groupoids}
\externaldocument[more-groupoids-]{more-groupoids}
\externaldocument[etale-]{etale}
\externaldocument[chow-]{chow}
\externaldocument[intersection-]{intersection}
\externaldocument[pic-]{pic}
\externaldocument[weil-]{weil}
\externaldocument[adequate-]{adequate}
\externaldocument[dualizing-]{dualizing}
\externaldocument[duality-]{duality}
\externaldocument[discriminant-]{discriminant}
\externaldocument[derham-]{derham}
\externaldocument[local-cohomology-]{local-cohomology}
\externaldocument[algebraization-]{algebraization}
\externaldocument[curves-]{curves}
\externaldocument[resolve-]{resolve}
\externaldocument[models-]{models}
\externaldocument[functors-]{functors}
\externaldocument[equiv-]{equiv}
\externaldocument[pione-]{pione}
\externaldocument[etale-cohomology-]{etale-cohomology}
\externaldocument[proetale-]{proetale}
\externaldocument[relative-cycles-]{relative-cycles}
\externaldocument[more-etale-]{more-etale}
\externaldocument[trace-]{trace}
\externaldocument[crystalline-]{crystalline}
\externaldocument[spaces-]{spaces}
\externaldocument[spaces-properties-]{spaces-properties}
\externaldocument[spaces-morphisms-]{spaces-morphisms}
\externaldocument[decent-spaces-]{decent-spaces}
\externaldocument[spaces-cohomology-]{spaces-cohomology}
\externaldocument[spaces-limits-]{spaces-limits}
\externaldocument[spaces-divisors-]{spaces-divisors}
\externaldocument[spaces-over-fields-]{spaces-over-fields}
\externaldocument[spaces-topologies-]{spaces-topologies}
\externaldocument[spaces-descent-]{spaces-descent}
\externaldocument[spaces-perfect-]{spaces-perfect}
\externaldocument[spaces-more-morphisms-]{spaces-more-morphisms}
\externaldocument[spaces-flat-]{spaces-flat}
\externaldocument[spaces-groupoids-]{spaces-groupoids}
\externaldocument[spaces-more-groupoids-]{spaces-more-groupoids}
\externaldocument[bootstrap-]{bootstrap}
\externaldocument[spaces-pushouts-]{spaces-pushouts}
\externaldocument[spaces-chow-]{spaces-chow}
\externaldocument[groupoids-quotients-]{groupoids-quotients}
\externaldocument[spaces-more-cohomology-]{spaces-more-cohomology}
\externaldocument[spaces-simplicial-]{spaces-simplicial}
\externaldocument[spaces-duality-]{spaces-duality}
\externaldocument[formal-spaces-]{formal-spaces}
\externaldocument[restricted-]{restricted}
\externaldocument[spaces-resolve-]{spaces-resolve}
\externaldocument[formal-defos-]{formal-defos}
\externaldocument[defos-]{defos}
\externaldocument[cotangent-]{cotangent}
\externaldocument[examples-defos-]{examples-defos}
\externaldocument[algebraic-]{algebraic}
\externaldocument[examples-stacks-]{examples-stacks}
\externaldocument[stacks-sheaves-]{stacks-sheaves}
\externaldocument[criteria-]{criteria}
\externaldocument[artin-]{artin}
\externaldocument[quot-]{quot}
\externaldocument[stacks-properties-]{stacks-properties}
\externaldocument[stacks-morphisms-]{stacks-morphisms}
\externaldocument[stacks-limits-]{stacks-limits}
\externaldocument[stacks-cohomology-]{stacks-cohomology}
\externaldocument[stacks-perfect-]{stacks-perfect}
\externaldocument[stacks-introduction-]{stacks-introduction}
\externaldocument[stacks-more-morphisms-]{stacks-more-morphisms}
\externaldocument[stacks-geometry-]{stacks-geometry}
\externaldocument[moduli-]{moduli}
\externaldocument[moduli-curves-]{moduli-curves}
\externaldocument[examples-]{examples}
\externaldocument[exercises-]{exercises}
\externaldocument[guide-]{guide}
\externaldocument[desirables-]{desirables}
\externaldocument[coding-]{coding}
\externaldocument[obsolete-]{obsolete}
\externaldocument[fdl-]{fdl}
\externaldocument[index-]{index}

% Heap Project documents
\externaldocument[linear-algebra-]{linear-algebra}
\externaldocument[groups-]{groups}
\externaldocument[complex-analysis-]{complex-analysis}
\externaldocument[functional-analysis-]{functional-analysis}
\externaldocument[categories-]{categories}
\externaldocument[commutative-algebra-]{commutative-algebra}
\externaldocument[ringed-spaces-]{ringed-spaces}
\externaldocument[homological-algebra-]{homological-algebra}
\externaldocument[lie-algebras-]{lie-algebras}
\externaldocument[representation-theory-]{representation-theory}
\externaldocument[smooth-manifolds-]{smooth-manifolds}
\externaldocument[differential-topology-]{differential-topology}
\externaldocument[algebraic-topology-]{algebraic-topology}
\externaldocument[differential-geometry-]{differential-geometry}
\externaldocument[algebraic-geometry-]{algebraic-geometry}
\externaldocument[symplectic-geometry-]{symplectic-geometry}
\externaldocument[complex-geometry-]{complex-geometry}
\externaldocument[hodge-theory-]{hodge-theory}
\maybeexternaldocument{deformations-}{deformations}
\externaldocument[vertex-algebras-]{vertex-algebras}
\externaldocument[springer-]{springer}
\externaldocument[infty-categories-]{infty-categories}
\externaldocument[seminars-differential-geometry-]{%
seminars-differential-geometry}
\externaldocument[seminars-complex-geometry-]{%
seminars-complex-geometry}
\externaldocument[seminars-algebraic-geometry-]{%
seminars-algebraic-geometry}
\externaldocument[seminars-others-]{%
seminars-others}
%\externaldocument[geometric-prequantization-]{geometric-prequantization}
\externaldocument[surfaces-]{surfaces}
\externaldocument[birational-maps-commutative-algebra-]{birational-maps-commutative-algebra}
\externaldocument[cremona-transformations-]{cremona-transformations}
\externaldocument[derived-categories-of-sheaves-]{derived-categories-of-sheaves}
\externaldocument[geometric-invariant-theory-]{geometric-invariant-theory}
\externaldocument[geometric-stability-conditions-and-group-actions-]{geometric-stability-conditions-and-group-actions}
\externaldocument[moduli-spaces-of-sheaves-]{moduli-spaces-of-sheaves}
\externaldocument[bridgeland-stability-general-theory-]{bridgeland-stability-general-theory}
\externaldocument[hilbert-schemes-]{hilbert-schemes}
\externaldocument[noncommutative-abelian-surfaces-and-kummer-type-hyperkahler-varieties-]{noncommutative-abelian-surfaces-and-kummer-type-hyperkahler-varieties}
\externaldocument[mumford-tate-groups-in-hodge-theory-]{mumford-tate-groups-in-hodge-theory}
\externaldocument[hodge-birational-atoms-2026-]{hodge-birational-atoms-2026}
\externaldocument[xxii-egd-2026-]{%
xxii-egd-2026}
\externaldocument[pde-]{pde}
\externaldocument[probability-]{probability}

\externaldocument[linguistics-]{linguistics}
\externaldocument[genealogia-familiar-]{%
genealogia-familiar}

\externaldocument[%
16th-alga-meeting-2026-]{%
16th-alga-meeting-2026}

\externaldocument[%
iv-jornadas-lefschetz-2026-]{%
iv-jornadas-lefschetz-2026}

\externaldocument[reading-the-masters-]{%
reading-the-masters}

\externaldocument[real-analysis-]{%
real-analysis}

\externaldocument[optimal-transport-]{%
optimal-transport}

\externaldocument[history-]{history}

% Theorem environments.
%
\theoremstyle{plain}
\newtheorem{theorem}[subsection]{Theorem}
\newtheorem{proposition}[subsection]{Proposition}
\newtheorem{lemma}[subsection]{Lemma}

\theoremstyle{definition}
\newtheorem{definition}[subsection]{Definition}
\newtheorem{example}[subsection]{Example}
\newtheorem{exercise}[subsection]{Exercise}
\newtheorem{situation}[subsection]{Situation}

\theoremstyle{remark}
\newtheorem{remark}[subsection]{Remark}
\newtheorem{remarks}[subsection]{Remarks}

\numberwithin{equation}{subsection}

% Macros
%
\def\lim{\mathop{\mathrm{lim}}\nolimits}
\def\colim{\mathop{\mathrm{colim}}\nolimits}
\def\Spec{\mathop{\mathrm{Spec}}}
\def\Hom{\mathop{\mathrm{Hom}}\nolimits}
\def\Ext{\mathop{\mathrm{Ext}}\nolimits}
\def\SheafHom{\mathop{\mathcal{H}\!\mathit{om}}\nolimits}
\def\SheafExt{\mathop{\mathcal{E}\!\mathit{xt}}\nolimits}
\def\Sch{\mathit{Sch}}
\def\Mor{\mathop{\mathrm{Mor}}\nolimits}
\def\Ob{\mathop{\mathrm{Ob}}\nolimits}
\def\Sh{\mathop{\mathit{Sh}}\nolimits}
\def\NL{\mathop{N\!L}\nolimits}
\def\CH{\mathop{\mathrm{CH}}\nolimits}
\def\proetale{{pro\text{-}\acute{e}tale}}
\def\etale{{\acute{e}tale}}
\def\QCoh{\mathit{QCoh}}
\def\Ker{\mathop{\mathrm{Ker}}}
\def\Im{\mathop{\mathrm{Im}}}
\def\Coker{\mathop{\mathrm{Coker}}}
\def\Coim{\mathop{\mathrm{Coim}}}

% Boxtimes
%
\DeclareMathSymbol{\boxtimes}{\mathbin}{AMSa}{"02}

%
% Macros for moduli stacks/spaces
%
\def\QCohstack{\mathcal{QC}\!\mathit{oh}}
\def\Cohstack{\mathcal{C}\!\mathit{oh}}
\def\Spacesstack{\mathcal{S}\!\mathit{paces}}
\def\Quotfunctor{\mathrm{Quot}}
\def\Hilbfunctor{\mathrm{Hilb}}
\def\Curvesstack{\mathcal{C}\!\mathit{urves}}
\def\Polarizedstack{\mathcal{P}\!\mathit{olarized}}
\def\Complexesstack{\mathcal{C}\!\mathit{omplexes}}
% \Pic is the operator that assigns to X its picard group, usage \Pic(X)
% \Picardstack_{X/B} denotes the Picard stack of X over B
% \Picardfunctor_{X/B} denotes the Picard functor of X over B
\def\Pic{\mathop{\mathrm{Pic}}\nolimits}
\def\Picardstack{\mathcal{P}\!\mathit{ic}}
\def\Picardfunctor{\mathrm{Pic}}
\def\Deformationcategory{\mathcal{D}\!\mathit{ef}}

\begin{document}

\title{IV Brazilian Workshop on
Lefschetz Properties}
\maketitle
\phantomsection
\label{section-phantom}
\tableofcontents

\medskip\noindent
UFBA, Salvador, Bahia.
September 28--October 2, 2026.

\href{%
https://jornadaslp.github.io/%
jornadas-lp/programacao.html%
}{Official programme.}
Talk titles and times checked on
September 28, 2026.
The sections follow the programme order.
Unannounced titles are marked explicitly.

\section{Introduction}
\label{section-introduction}

\noindent
These notes start from Daniel's notes of
the introductory discussion on
September 28, 2026.
The conventions and qualifications below
make the geometric motivation precise.
The HP starting points are
\emph{Complex Geometry},
``Lefschetz theorems'', and
\emph{Commutative Algebra},
``Artinian algebras''.

\begin{theorem}[Hard Lefschetz]
\label{theorem-hard-lefschetz}
Let $X$ be a compact Kähler manifold of
complex dimension $n$, with Kähler form
$\omega$.
Define the Lefschetz operator by
$$
\begin{aligned}
L:H^j(X,\mathbb{R})
&\longrightarrow H^{j+2}(X,\mathbb{R}),\\
\alpha&\longmapsto[\omega]\smile\alpha.
\end{aligned}
$$
For every $0\leq j\leq n$, the map
$$
L^{n-j}:H^j(X,\mathbb{R})
\longrightarrow H^{2n-j}(X,\mathbb{R})
$$
is an isomorphism.
For a smooth complex projective variety,
one may take the first Chern class of an
ample line bundle and use rational
coefficients.
\end{theorem}

\medskip\noindent
\textbf{Motivational question.}
Quais propriedades deve satisfazer uma
álgebra para ser candidata a ser a
cohomologia de uma variedade?

For a connected smooth complex projective
variety, its rational cohomology is
finite-dimensional and graded-commutative,
with a unit and $H^0=\mathbb{Q}$.
It satisfies Poincaré duality and Hard
Lefschetz for an ample class, and carries
further Hodge-theoretic structure.
These are necessary conditions; we are not
claiming that they are sufficient.

\begin{remark}[Which grading?]
\label{remark-grading}
Multiplication by a Kähler class raises
cohomological degree by two.
To compare with commutative algebra, put
$$
A_k=H^{2k}(X,\mathbb{Q}).
$$
The ample class now has degree one.
The even cohomology is commutative, but
need not be generated by $A_1$.
Moreover, odd cohomology can be nonzero:
the entire cohomology algebra is generally
only graded-commutative.
Thus a standard graded commutative algebra
is a more restrictive model than the
cohomology algebra of an arbitrary variety.
\end{remark}

\begin{example}[Products of projective
spaces]
\label{example-projective-products}
Let $a_0,\ldots,a_r$ be positive integers
and put
$$
X=\prod_{i=0}^{r}\mathbb{P}^{a_i-1}.
$$
If $x_i$ denotes the pullback of the
hyperplane class from the $i$th factor,
then
$$
H^*(X,\mathbb{Q})\cong
\frac{\mathbb{Q}[x_0,\ldots,x_r]}
{(x_0^{a_0},\ldots,x_r^{a_r})}.
$$
Each $x_i$ has cohomological degree two;
after regrading, it has algebra degree one.
The class
$$
\ell=x_0+\cdots+x_r
$$
is ample, and its Lefschetz operator is
multiplication by $\ell$.
Writing $N=\sum_{i=0}^{r}(a_i-1)$, Hard
Lefschetz gives isomorphisms
$$
\times\ell^{N-2k}:A_k\longrightarrow A_{N-k}
$$
for $0\leq k\leq N/2$.
The Lefschetz decomposition then gives
maximal rank for every power between any
two graded pieces.
\end{example}

For instance, for
$X=\mathbb{P}^1\times\mathbb{P}^1$,
$$
A=\mathbb{Q}[x,y]/(x^2,y^2),
\qquad \ell=x+y.
$$
The graded dimensions are $(1,2,1)$.
Multiplication sends $1$ to $x+y$,
sends both $x$ and $y$ to $xy$, and
$$
\ell^2=2xy.
$$
Thus $A_0\longrightarrow A_1$ is
injective, $A_1\longrightarrow A_2$ is
surjective, and
$\times\ell^2:A_0\longrightarrow A_2$
is an isomorphism.

\medskip\noindent
\textbf{Realization question.}
Which graded algebras can be realized as
the rational cohomology algebra of a
smooth complex projective variety?
This is a broad realization problem,
not a criterion supplied by SLP alone.
One must specify the class of varieties,
coefficients, and grading when asking it.

\begin{definition}[Standard graded algebra]
\label{definition-standard-graded}
Let $K$ be a field of characteristic zero.
A \emph{standard graded algebra} is a
commutative graded $K$-algebra
$$
A=\bigoplus_{k\geq0}A_k
$$
with $A_0=K$, generated by a
finite-dimensional space $A_1$.
Equivalently, $A=Q/I$ for a polynomial ring
$Q=K[X_0,\ldots,X_r]$, with variables of
degree one and $I$ homogeneous.
Here an \emph{Artinian algebra} means that
$A$ is finite-dimensional over $K$, as in
the HP notes.
For standard graded algebras this agrees
with the descending-chain definition.
\end{definition}

\begin{definition}[Maximal rank]
\label{definition-maximal-rank}
For $k\geq0$ and $t\geq1$, say that $A$
has the \emph{maximal rank property in
degree $k$ and range $t$} if there is a
linear form $\ell\in A_1$ such that
$$
\times\ell^t:A_k\longrightarrow A_{k+t}
$$
has rank
$$
\min\{\dim_K A_k,\dim_K A_{k+t}\}.
$$
Thus the map is injective or surjective,
as dictated by the two dimensions.
\end{definition}

\begin{remark}
\label{remark-maximal-rank-terminology}
The target is $A_{k+t}$ and the linear
form belongs to $A_1$.
This is the local terminology of these
notes. In the literature, ``Maximal Rank
Property'' can instead refer to
multiplication by general forms of each
degree, which need not be powers of a
linear form.
\end{remark}

\begin{definition}[Strong Lefschetz
property]
\label{definition-slp}
An Artinian standard graded algebra $A$
has the \emph{Strong Lefschetz Property}
(SLP) if there is a single
$\ell\in A_1$ for which
$$
\times\ell^t:A_k\longrightarrow A_{k+t}
$$
has maximal rank for every $k\geq0$ and
every $t\geq1$.
Such an $\ell$ is a \emph{strong Lefschetz
element}.
\end{definition}

\medskip\noindent
\textbf{Question.}
Can we characterize all Artinian standard
graded algebras having SLP?

\begin{remark}[Why Artinian algebras?]
\label{remark-why-artinian}
Estudamos álgebras artinianas pensando
em anéis de cohomologia de variedades.

Daniel's question: então toda álgebra de
cohomologia é artiniana?
For a compact smooth manifold, the total
cohomology over a field is
finite-dimensional.
The same finiteness holds for complex
algebraic varieties of finite type.
So the finite-dimensional part of the
intuition is correct in these settings.
It is not true for arbitrary topological
spaces.

There is still a commutativity qualification:
the full cohomology is graded-commutative.
Its even part is a commutative
finite-dimensional algebra, hence Artinian
in the HP sense.
When odd cohomology vanishes, this
qualification disappears for the full ring.
Standard generation in degree one after
regrading is an additional condition.
\end{remark}

\begin{definition}[Weak Lefschetz property]
\label{definition-wlp}
An Artinian standard graded algebra $A$
has the \emph{Weak Lefschetz Property}
(WLP) if there is a single
$\ell\in A_1$ such that
$$
\times\ell:A_k\longrightarrow A_{k+1}
$$
has maximal rank for every $k\geq0$.
Such an $\ell$ is a \emph{weak Lefschetz
element}.
\end{definition}

SLP implies WLP, by taking $t=1$.
The converse does not hold in general.

\begin{lemma}[Change of variables]
\label{lemma-change-of-variables}
SLP and WLP are invariant under graded
algebra isomorphisms, in particular under
invertible linear changes of variables.
\end{lemma}

\begin{proof}
If $\varphi:A\longrightarrow B$ is a
graded algebra isomorphism, then
$$
\varphi(\ell^t a)=\varphi(\ell)^t\varphi(a).
$$
The isomorphisms on the graded pieces
therefore identify the multiplication maps
and preserve their ranks.
\end{proof}

\begin{lemma}[Powers of independent forms]
\label{lemma-powers-of-forms}
Let $\ell_0,\ldots,\ell_r$ be linearly
independent linear forms in
$Q=K[X_0,\ldots,X_r]$.
For positive integers $a_i$, the algebra
$$
Q/(\ell_0^{a_0},\ldots,\ell_r^{a_r})
$$
has SLP in characteristic zero.
\end{lemma}

\begin{proof}
The invertible substitution
$X_i\longmapsto\ell_i$ identifies this
algebra with the monomial complete
intersection
$$
Q/(X_0^{a_0},\ldots,X_r^{a_r}).
$$
Over $\mathbb{Q}$, its SLP follows from the
product-of-projective-spaces example.
The multiplication matrices for the element
$X_0+\cdots+X_r$ have integer entries.
Their nonzero maximal minors stay nonzero
over every field of characteristic zero.
Apply the change-of-variables lemma.
\end{proof}

Linear independence is essential for this
argument: arbitrary linear forms need not
define a change of variables or even give
an Artinian quotient.

\begin{definition}[Artinian complete
intersection]
\label{definition-complete-intersection}
An \emph{Artinian complete intersection}
is an algebra
$$
A=K[X_0,\ldots,X_r]/(f_0,\ldots,f_r),
$$
where the homogeneous forms $f_i$ form a
regular sequence and the quotient is
finite-dimensional.
A \emph{regular sequence} means that each
$f_i$ is a non-zero-divisor modulo the
preceding forms, and the generated ideal
is proper.
The algebra is \emph{equigenerated} if all
the $f_i$ have the same degree.
If there are no linear relations, its
codimension is $r+1=\dim_K A_1$.
\end{definition}

\begin{lemma}[Generic complete
intersections]
\label{lemma-generic-complete-intersections}
Fix positive degrees $a_0,\ldots,a_r$.
A general Artinian complete intersection
of those degrees has SLP in
characteristic zero.
\end{lemma}

\begin{proof}
The complete-intersection locus in the
parameter space of tuples of forms of
these degrees is a nonempty Zariski-open
set. Its quotients have a fixed Hilbert
function and vanish above degree
$\sum_i(a_i-1)$.
Locally in this parameter space, choose
bases for the graded pieces.
For a fixed linear form, maximal rank of
each multiplication map is an open
condition: some maximal minor is nonzero.
Only finitely many maps are relevant.

At the tuple $(X_0^{a_0},\ldots,X_r^{a_r})$,
the form $X_0+\cdots+X_r$ works for all
these maps. Their simultaneous maximal
rank therefore holds on a nonempty open
set of tuples.
\end{proof}

Here ``general'' or ``generic'' means
belonging to a suitable nonempty
Zariski-open set in the parameter space
with the degrees fixed.
This conclusion uses openness in addition
to invariance under change of variables.
It does not say that every complete
intersection is a change of variables of
a monomial one.

\medskip\noindent
\textbf{The big problem.}
Does every Artinian complete intersection
in characteristic zero have WLP?
Does every such algebra have SLP?
The conjectural answer to both is yes.
The following established cases concern
\emph{WLP}, and must not be quoted as
general SLP theorems:

\begin{enumerate}
\item Every codimension-three Artinian
complete intersection has WLP.
This is the theorem of Harima, Migliore,
Nagel and Watanabe; see
\href{%
https://arxiv.org/abs/math/0208201%
}{The Weak and Strong Lefschetz
Properties for Artinian K-Algebras}.

\item In codimension four, equigenerated
complete intersections with common
relation degree $2\leq d\leq5$ have WLP.
Boij, Migliore, Miró-Roig and Nagel prove
the cases $d=3,4,5$ and recall the known
quadratic case in their introduction;
see \href{%
https://arxiv.org/abs/2212.09890%
}{On the Weak Lefschetz Property for
height four equigenerated complete
intersections}.
Here $d$ is the degree of each defining
form, not the socle degree.
\end{enumerate}

\begin{remark}[Characteristic matters]
\label{remark-characteristic}
The characteristic-zero hypothesis is
part of these statements.
For example, in characteristic two,
every linear form in $K[x,y]/(x^2,y^2)$
has square zero, so this algebra fails
SLP although it is a monomial complete
intersection.
\end{remark}

\section{Compressed Artinian Gorenstein
algebras}
\label{section-lenin-bezerra}

Lenin Alexandre de Almeida Bezerra (UPE).

Segunda 28/09, 11:40.

\medskip\noindent
\textbf{Notes.}
To be added during the talk.

\section{Expect the unexpected: probabilistic
failure of WLP in monomial algebras}
\label{section-thiago-holleben}

Thiago Holleben (Max Planck Institute).

Segunda 28/09, 14:30.

\medskip\noindent
\textbf{Notes.}
To be added during the talk.

\section{Milnor Algebras, Higher Jacobians
and
Lefschetz-Type Maps}
\label{section-giovanna-ilardi}

Giovanna Ilardi (Università di Napoli).

Terça 29/09, 9:30.

\medskip\noindent
\textbf{Notes.}
To be added during the talk.

\section{Homaloidal hypersurfaces and a
product
formula for gradient maps}
\label{section-nivaldo-junior}

Nivaldo Nunes de Medeiros Júnior (UFF).

Terça 29/09, 10:30.

\medskip\noindent
\textbf{Notes.}
To be added during the talk.

\section{Title to be announced: Kézia
Patrícia
Mestre Carvalho}
\label{section-kezia-carvalho}

Kézia Patrícia Mestre Carvalho (UFPE).

Terça 29/09, 11:40.

The official programme has not yet
announced a title.

\medskip\noindent
\textbf{Notes.}
To be added during the talk.

\section{Title to be announced: João Paulo
Costalonga}
\label{section-joao-costalonga}

João Paulo Costalonga (UFES).

Terça 29/09, 14:30.

The official programme has not yet
announced a title.

\medskip\noindent
\textbf{Notes.}
To be added during the talk.

\section{Bourbaki schemes of three generated
ideals}
\label{section-marcos-jardim}

Marcos Jardim (UNICAMP).

Quarta 30/09, 9:30.

\medskip\noindent
\textbf{Notes.}
To be added during the talk.

\section{Non-dominant quadratic maps and
cubics with
vanishing Hessian}
\label{section-thiago-fassarela}

Thiago Fassarela (UFF).

Quarta 30/09, 10:30.

\medskip\noindent
\textbf{Notes.}
To be added during the talk.

\section{On Artinian Gorenstein Algebras with
Macaulay Dual Generator of Fixed Waring
Rank}
\label{section-janaine-martins}

Janaíne Geralda Mesquita Martins (UFMG).

Quarta 30/09, 11:40.

\medskip\noindent
\textbf{Notes.}
To be added during the talk.

\section{Higgs Grassmannians}
\label{section-ugo-bruzzo}

Ugo Bruzzo (UFMG).

Quinta 01/10, 9:30.

\medskip\noindent
\textbf{Notes.}
To be added during the talk.

\section{Almost Cohen-Macaulay height two
ideals
equigenerated by three forms}
\label{section-thiago-cabral}

Thiago Fiel da Costa Cabral (UFPE).

Quinta 01/10, 10:30.

\medskip\noindent
\textbf{Notes.}
To be added during the talk.

\section{Multiple Tor modules}
\label{section-rafael-holanda}

Rafael Ferreira Holanda (UFPE).

Quinta 01/10, 11:40.

\medskip\noindent
\textbf{Notes.}
To be added during the talk.

\section{Hessians, curvature and lightlike
hypersurfaces}
\label{section-thiago-silva}

Thiago Dias de Oliveira Silva (UFRPE).

Quinta 01/10, 14:30.

\medskip\noindent
\textbf{Notes.}
To be added during the talk.

\section{Veronese-Avoiding Hypersurfaces and
Lefschetz Properties}
\label{section-abbas-nejad}

Abbas Nasrollah Nejad (UNIFESP).

Sexta 02/10, 9:30.

\medskip\noindent
\textbf{Notes.}
To be added during the talk.

\section{Newton Duality and Macaulay Inverse
System
of Gorenstein Ideals}
\label{section-dayane-lira}

Dayane Santos de Lira (UFERSA).

Sexta 02/10, 10:30.

\medskip\noindent
\textbf{Notes.}
To be added during the talk.

\end{document}
